\appendix

\section{Proof Details}
\setlength{\emergencystretch}{3em}

\subsection{Root Feasibility and Finite Pivot Conditioning}

\begin{lemma}[Compact feasibility of cube-supported roots]
\label{lem:root-feasible-compact}
Under Assumption~\ref{assump:shared-pfaffian-chain}, for $N\geq1$, $R>0$,
and compact interval $\Theta$, the set
\[
K_R=\left\{\theta\in\Theta:
 |b(\theta)|\leq R\sum_{i=1}^N|F_i(\theta)|\right\}
\]
is compact.  For every $\alpha\in[-R,R]^N$ and $\theta\in\Theta$,
$\phi_\alpha(\theta)=0$ implies $\theta\in K_R$.  If $K_R=\varnothing$,
then
\[
\{\alpha\in[-R,R]^N:\exists\theta\in\Theta,
 \ \phi_\alpha(\theta)=0\}=\varnothing,
\qquad
\Gamma_{\rm piv}(b,F;R)=0.
\]
\end{lemma}

\begin{proof}
Assumption~\ref{assump:shared-pfaffian-chain} makes $b,F_1,\ldots,F_N$
continuous.  Hence
\[
\theta\longmapsto |b(\theta)|-R\sum_{i=1}^N|F_i(\theta)|
\]
is continuous, and $K_R$ is its inverse image of $(-\infty,0]$.  It is
closed relative to compact $\Theta$, and therefore compact.  This relative
topology argument retains both endpoints of $\Theta$.

If $\alpha\in[-R,R]^N$ and $\phi_\alpha(\theta)=0$, then
$b(\theta)=-\sum_i\alpha_iF_i(\theta)$, so the finite triangle inequality
gives
\[
\begin{aligned}
|b(\theta)|
&=\left|\sum_{i=1}^N\alpha_iF_i(\theta)\right|\\
&\leq\sum_{i=1}^N|\alpha_i|\,|F_i(\theta)|\\
&\leq R\sum_{i=1}^N|F_i(\theta)|
=R\|F(\theta)\|_1.
\end{aligned}
\]
Thus $\theta\in K_R$.  When $K_R$ is empty, this implication rules out
every cube-supported root.  The second branch in the definition of
$\Gamma_{\rm piv}$ then gives zero; no empty-supremum convention is being
used here.
\end{proof}

\begin{lemma}[Positive available-pivot margin]
\label{lem:positive-pivot-margin}
Under Assumptions~\ref{assump:shared-pfaffian-chain} and
\ref{assump:no-forced-root} and Lemma~\ref{lem:root-feasible-compact}, if
$K_R\neq\varnothing$, then $F(\theta)\neq0$ for every $\theta\in K_R$ and
\[
\rho:=\min_{\theta\in K_R}\max_{1\leq j\leq N}|F_j(\theta)|>0.
\]
At each $\theta\in K_R$, every index attaining the displayed coordinate
maximum is a legal pivot and has $|F_j(\theta)|\geq\rho$.
\end{lemma}

\begin{proof}
If $\theta\in K_R$ and $F(\theta)=0$, feasibility gives
\[
|b(\theta)|\leq R\|F(\theta)\|_1=0.
\]
Then $(b(\theta),F(\theta))=(0,0)$, contrary to
Assumption~\ref{assump:no-forced-root}.  This argument is confined to
$K_R$; a point outside $K_R$ with $F=0$ and $b\neq0$ remains admissible and
root-free.

The function
\[
g(\theta)=\max_{1\leq j\leq N}|F_j(\theta)|
\]
is continuous as a finite maximum of continuous functions, and it is
strictly positive on $K_R$.  Lemma~\ref{lem:root-feasible-compact} makes
$K_R$ nonempty and compact in the present branch, so $g$ attains a minimum
at some $\theta_*\in K_R$.  That minimum is
$\rho=g(\theta_*)>0$.  The finite coordinate set has a maximizer at each
point, and every maximizer has magnitude $g(\theta)\geq\rho$.

Coordinates may vanish and maximizers may tie; at least one maximizer is
nonzero.  The choice here is pointwise, with no assertion of a continuous
or measurable global index.  When $N=1$, the sole coordinate is the
maximizer and is nonzero on $K_R$.  Compact-set endpoints are included and
no transversality condition is used.
\end{proof}

\begin{proposition}[Crude fixed-family pivot-variation bound]
\label{prop:fixed-family-pivot-bound}
Under Assumptions~\ref{assump:shared-pfaffian-chain} and
\ref{assump:no-forced-root} and Lemmas~\ref{lem:root-feasible-compact} and
\ref{lem:positive-pivot-margin}, suppose $K_R\neq\varnothing$ and define
\[
B_0:=\max_{\theta\in\Theta}
\max\{|b(\theta)|,|F_1(\theta)|,\ldots,|F_N(\theta)|\},
\]
\[
B_1:=\max_{\theta\in\Theta}
\max\{|b'(\theta)|,|F_1'(\theta)|,\ldots,|F_N'(\theta)|\}.
\]
Then $B_0,B_1<\infty$, and for every $\theta\in K_R$,
\[
\min_{1\leq k\leq N}V_k(\theta)
\leq2B_0B_1\rho^{-2}\bigl(1+R(N-1)\bigr).
\]
Consequently,
\[
\Gamma_{\rm piv}(b,F;R)
\leq2B_0B_1\rho^{-2}\bigl(1+R(N-1)\bigr)<\infty.
\]
\end{proposition}

\begin{proof}
The functions $b,F_i,b',F_i'$ are continuous on compact $\Theta$, so the
two displayed maxima exist and are finite.  Nonempty feasibility also gives
$B_0\geq\rho>0$, although we never divide by $B_0$; $B_1$ need not be
positive.

Fix $\theta\in K_R$ and choose an index $j$ attaining
$\max_k|F_k(\theta)|$.  Lemma~\ref{lem:positive-pivot-margin} gives
$|F_j(\theta)|\geq\rho>0$.  The quotient rule and the two-term triangle
inequality give
\[
\begin{aligned}
\left|\left(\frac b{F_j}\right)'(\theta)\right|
&=\frac{|b'(\theta)F_j(\theta)-b(\theta)F_j'(\theta)|}
{|F_j(\theta)|^2}\\
&\leq\frac{|b'(\theta)|\,|F_j(\theta)|
 +|b(\theta)|\,|F_j'(\theta)|}{|F_j(\theta)|^2}\\
&\leq\frac{B_1B_0+B_0B_1}{\rho^2}
=\frac{2B_0B_1}{\rho^2}.
\end{aligned}
\]
For each $i\neq j$, the separate quotient calculation is
\[
\begin{aligned}
\left|\left(\frac{F_i}{F_j}\right)'(\theta)\right|
&=\frac{|F_i'(\theta)F_j(\theta)-F_i(\theta)F_j'(\theta)|}
{|F_j(\theta)|^2}\\
&\leq\frac{|F_i'(\theta)|\,|F_j(\theta)|
 +|F_i(\theta)|\,|F_j'(\theta)|}{|F_j(\theta)|^2}\\
&\leq\frac{2B_0B_1}{\rho^2}.
\end{aligned}
\]
There is one offset ratio and exactly $N-1$ feature ratios.  Therefore
\[
\begin{aligned}
V_j(\theta)
&=\left|\left(\frac b{F_j}\right)'(\theta)\right|
 +R\sum_{i\neq j}\left|\left(\frac{F_i}{F_j}\right)'(\theta)\right|\\
&\leq\frac{2B_0B_1}{\rho^2}
 +R(N-1)\frac{2B_0B_1}{\rho^2}\\
&=2B_0B_1\rho^{-2}\bigl(1+R(N-1)\bigr).
\end{aligned}
\]
Although zero pivots have extended speed $+\infty$, the selected $j$ is
nonzero; hence $\min_kV_k(\theta)\leq V_j(\theta)<\infty$.  Taking the
supremum over nonempty $K_R$ proves the result.

If an unselected $F_i(\theta)$ vanishes, its numerator term
$F_iF_j'$ is zero while $F_i'F_j$ is covered by the displayed bound.  If
$b(\theta)=0$, its corresponding term vanishes; if $b\equiv0$, then
$b'\equiv0$ and the offset-ratio derivative is zero.  When $B_1=0$, every
displayed derivative is zero.  For $N=1$ the feature-ratio sum is empty and
$1+R(N-1)=1$.  The $C^1$ convention includes endpoints.  This certificate
is fixed-family: it gives no bound for $B_0,B_1$, or $\rho^{-1}$ in terms of
$q,M,\Delta_{\rm rnd},\Delta_{\rm aff}$ or other general presentation data.
\end{proof}

\begin{proposition}[Root feasibility and finite pivot conditioning]
\label{prop:root-feasibility-finite-conditioning}
Under Assumptions~\ref{assump:shared-pfaffian-chain} and
\ref{assump:no-forced-root}, every root generated by
$\alpha\in[-R,R]^N$ lies in compact $K_R$, a legal pivot is available at
every feasible point, and
\[
0\leq\Gamma_{\rm piv}(b,F;R)<\infty.
\]
If $K_R=\varnothing$, the supported-root event is empty and
$\Gamma_{\rm piv}=0$; otherwise the bound in
Proposition~\ref{prop:fixed-family-pivot-bound} holds.
\end{proposition}

\begin{proof}
Lemma~\ref{lem:root-feasible-compact} gives compactness, the supported-root
implication, and the complete empty branch.  In the nonempty branch,
Lemma~\ref{lem:positive-pivot-margin} gives $\rho>0$ and a nonzero
largest-coordinate pivot at each point.  Proposition~\ref{prop:fixed-family-pivot-bound}
bounds all quotient terms for such a pivot and then takes the defining
supremum.  These two branches prove every assertion.  The quantities
$B_0,B_1,\rho$ witness finiteness but are not additional public rate
factors.
\end{proof}

\subsection{Measurable Adaptive Coordinate Charts}

\begin{lemma}[Borel measurability of extended pivot speeds]
\label{lem:extended-speed-borel}
Under Assumption~\ref{assump:shared-pfaffian-chain} and
Proposition~\ref{prop:root-feasibility-finite-conditioning}, every extension
$V_j:\Theta\to[0,+\infty]$ is Borel.  If $K_R\neq\varnothing$, then for
every $\theta\in K_R$,
\[
\min_{1\leq i\leq N}V_i(\theta)
\leq\Gamma_{\rm piv}(b,F;R)<\infty.
\]
For empty $K_R$ the pointwise statement is vacuous and
$\Gamma_{\rm piv}=0$.
\end{lemma}

\begin{proof}
Fix $j$.  Continuity of $F_j$ makes
$U_j=\{\theta:F_j(\theta)\neq0\}$ relatively open and Borel.  On $U_j$,
\[
\left(\frac b{F_j}\right)'=\frac{b'F_j-bF_j'}{F_j^2},
\qquad
\left(\frac{F_i}{F_j}\right)'=
\frac{F_i'F_j-F_iF_j'}{F_j^2}\quad(i\neq j).
\]
All numerator factors are continuous and the denominator is nonzero there,
so
\[
W_j(\theta)=
\left|\left(\frac b{F_j}\right)'(\theta)\right|
 +R\sum_{i\neq j}\left|\left(\frac{F_i}{F_j}\right)'(\theta)\right|
\]
is finite and continuous on $U_j$, including in the relative topology at
endpoints.  Since $V_j=W_j$ on $U_j$ and $V_j=+\infty$ elsewhere, for every
finite $a$,
\[
\{\theta:V_j(\theta)<a\}=\{\theta\in U_j:W_j(\theta)<a\},
\]
which is Borel.  Finite strict sublevel sets generate the Borel structure on
$[0,+\infty]$, so $V_j$ is Borel.  The same argument proves
\[
V_j(\theta)<+\infty\quad\Longleftrightarrow\quad F_j(\theta)\neq0.
\]

For nonempty $K_R$, Proposition~\ref{prop:root-feasibility-finite-conditioning}
provides a finite ordinary speed at each feasible point.  A finite list with
one finite entry has a finite attained minimum, and the definition of
$\Gamma_{\rm piv}$ gives the displayed bound.  The empty branch follows
from the same proposition.  No continuity across a zero of $F_j$ is
asserted.
\end{proof}

\begin{proposition}[Borel lexicographic pivot partition]
\label{prop:lexicographic-pivot-partition}
Under Lemma~\ref{lem:extended-speed-borel} and
Lemma~\ref{lem:root-feasible-compact}, define for $\theta\in K_R$
\[
j_*(\theta)=\min\operatorname*{arg\,min}_{1\leq i\leq N}V_i(\theta).
\]
Then $j_*:K_R\to\{1,\ldots,N\}$ is Borel.  Its fibers
$E_j=\{\theta\in K_R:j_*(\theta)=j\}$ obey
\[
E_j=K_R\cap\{V_j<+\infty\}
\cap\bigcap_{i<j}\{V_j<V_i\}
\cap\bigcap_{i>j}\{V_j\leq V_i\}.
\tag{A1}\label{eq:a1}
\]
They are Borel, pairwise disjoint, and cover $K_R$.  On $E_j$,
\[
V_j(\theta)=\min_iV_i(\theta)<+\infty,
\qquad F_j(\theta)\neq0.
\]
These conclusions also hold for empty $K_R$, with empty fibers.
\end{proposition}

\begin{proof}
For extended-real Borel functions $f,g$,
\[
\{f<g\}=\bigcup_{r\in\mathbb Q}
\bigl(\{f<r\}\cap\{g>r\}\bigr),
\qquad
\{f\leq g\}=\Theta\setminus\{g<f\}.
\tag{A2}\label{eq:a2}
\]
The first equality follows by rational separation when $g$ is finite and,
when $g=+\infty$, by choosing a finite rational larger than $f$.  Thus all
comparisons in \eqref{eq:a1} are Borel, as are $K_R$ and
$\{V_j<+\infty\}$.

At each $\theta\in K_R$, Lemma~\ref{lem:extended-speed-borel} gives a
finite attained minimum.  If $j$ is its least attaining index, then
$V_j<V_i$ for $i<j$ and $V_j\leq V_i$ for $i>j$, placing the point in
$E_j$.  Conversely, \eqref{eq:a1} makes $j$ the least minimizer.  This
proves coverage and uniqueness, hence disjointness.  Each singleton fiber
of $j_*$ is Borel; inverse images of subsets of the finite discrete target
are finite unions of these fibers, so $j_*$ is Borel.  Finiteness of $V_j$
on its fiber and Lemma~\ref{lem:extended-speed-borel} give $F_j\neq0$.

When $N=1$, both comparison intersections are empty, $E_1=K_R$, and the
selector is one on its domain.  Empty $K_R$ gives the unique empty map.
Zeros in unselected coordinates are unrestricted.
\end{proof}

\begin{lemma}[Exact denominator exhaustion]
\label{lem:denominator-exhaustion}
Under Assumption~\ref{assump:shared-pfaffian-chain} and
Proposition~\ref{prop:lexicographic-pivot-partition}, for $m\geq1$ define
\[
E_{j,m}=E_j\cap\{\theta\in\Theta:|F_j(\theta)|\geq1/m\}.
\]
Then every $E_{j,m}$ is Borel and
\[
E_{j,m}\subseteq E_{j,m+1},
\qquad
\bigcup_{m=1}^\infty E_{j,m}=E_j.
\tag{A3}\label{eq:a3}
\]
Every selected point enters at a finite level.
\end{lemma}

\begin{proof}
Continuity of $F_j$ makes $\{|F_j|\geq1/m\}$ relatively closed, hence
Borel; intersecting with Borel $E_j$ preserves Borelness.  Since
$1/(m+1)\leq1/m$, the sets increase.  Their union is contained in $E_j$.
For the reverse inclusion, if $\theta\in E_j$ then
$x=|F_j(\theta)|>0$ by Proposition~\ref{prop:lexicographic-pivot-partition}.
There is a finite integer $m$ with $1/m\leq x$, so
$\theta\in E_{j,m}$.  The weak threshold retains equality points.  An
arbitrarily small positive pivot may enter late but leaves no remainder;
empty cells satisfy the same identity.
\end{proof}

\begin{proposition}[Root chart in the original coefficient coordinates]
\label{prop:original-coordinate-chart}
Under Assumption~\ref{assump:shared-pfaffian-chain} and
Proposition~\ref{prop:lexicographic-pivot-partition}, fix $j$ and set
$J_j=\{1,\ldots,N\}\setminus\{j\}$, ordered increasingly.  For
$\beta=(\beta_i)_{i\in J_j}$ define
\[
T_j(\theta,\beta)
=-\frac{b(\theta)+\sum_{i\in J_j}\beta_iF_i(\theta)}{F_j(\theta)}.
\tag{A4}\label{eq:a4}
\]
This function is well-defined and jointly Borel on
$E_j\times[-R,R]^{J_j}$.  If
\[
\alpha_i^{(j)}(\theta,\beta)=
\begin{cases}
T_j(\theta,\beta),&i=j,\\
\beta_i,&i\in J_j,
\end{cases}
\]
then
\[
b(\theta)+\langle\alpha^{(j)}(\theta,\beta),F(\theta)\rangle=0.
\tag{A5}\label{eq:a5}
\]
The conclusion does not assert $T_j\in[-R,R]$.
\end{proposition}

\begin{proof}
The selected cell lies in $U_j$.  On
$U_j\times\mathbb R^{J_j}$, rewrite \eqref{eq:a4} as
\[
-\frac{b(\theta)}{F_j(\theta)}
-\sum_{i\in J_j}\beta_i\frac{F_i(\theta)}{F_j(\theta)}.
\]
The coefficient ratios are continuous on $U_j$ and the expression is
finite affine in $\beta$, hence jointly continuous.  Its restriction to
the stated Borel domain is Borel.  Direct substitution gives
\[
\begin{aligned}
b+\langle\alpha^{(j)},F\rangle
&=b+T_jF_j+\sum_{i\in J_j}\beta_iF_i\\
&=b-\left(b+\sum_{i\in J_j}\beta_iF_i\right)
 +\sum_{i\in J_j}\beta_iF_i=0.
\end{aligned}
\]
Thus the residual vanishes in the original coefficient vector; there is no
weighted or transformed target.  Unselected features may vanish.  If
$N=1$, $J_1=\varnothing$, the beta space is the singleton
$\mathbb R^0$, both sums are empty, and \eqref{eq:a5} reads
$b+(-b/F_1)F_1=0$.  Empty cells give empty domains, and the relative $C^1$
convention includes interval endpoints.
\end{proof}

\begin{proposition}[Selected-chart velocity]
\label{prop:selected-chart-velocity}
Under Assumption~\ref{assump:shared-pfaffian-chain},
Proposition~\ref{prop:root-feasibility-finite-conditioning},
Proposition~\ref{prop:lexicographic-pivot-partition}, and
Proposition~\ref{prop:original-coordinate-chart}, the following holds for
every $\theta\in E_j$ and every $\beta\in[-R,R]^{J_j}$:
\[
\partial_\theta T_j(\theta,\beta)
=-\left(\frac b{F_j}\right)'(\theta)
-\sum_{i\in J_j}\beta_i
 \left(\frac{F_i}{F_j}\right)'(\theta),
\tag{A6}\label{eq:a6}
\]
and
\[
|\partial_\theta T_j(\theta,\beta)|
\leq V_j(\theta)
=\min_{1\leq i\leq N}V_i(\theta)
\leq\Gamma_{\rm piv}(b,F;R)<\infty.
\tag{A7}\label{eq:a7}
\]
The same conclusions hold on every exhausted domain
$E_{j,m}\times[-R,R]^{J_j}$.
\end{proposition}

\begin{proof}
The selected denominator is nonzero.  On $U_j$, the quotient identities
used in Lemma~\ref{lem:extended-speed-borel} hold.  Differentiating the
finite affine expression \eqref{eq:a4} at fixed $\beta$ gives
\eqref{eq:a6} without interchanging any limit or infinite sum.  At an
endpoint this is the derivative in the stated $C^1$ convention.

Using $|\beta_i|\leq R$ separately for every nonpivot coordinate,
\[
\begin{aligned}
|\partial_\theta T_j|
&\leq\left|\left(\frac b{F_j}\right)'\right|
 +\sum_{i\in J_j}|\beta_i|
  \left|\left(\frac{F_i}{F_j}\right)'\right|\\
&\leq\left|\left(\frac b{F_j}\right)'\right|
 +R\sum_{i\in J_j}
  \left|\left(\frac{F_i}{F_j}\right)'\right|=V_j.
\end{aligned}
\]
Cell membership gives $V_j=\min_iV_i$, and
$E_j\subseteq K_R$ together with
Proposition~\ref{prop:root-feasibility-finite-conditioning} gives the last
bound in \eqref{eq:a7}.  There is one literal factor $R$ per nonpivot term
and no sum over chart indices.  Restriction to $E_{j,m}$ changes no formula
and introduces no $m$-dependent constant.  For $N=1$ the sum is empty and
$|-(b/F_1)'|=V_1$.  If $K_R$ is empty, the domain is empty and
$\Gamma_{\rm piv}=0$.
\end{proof}

\begin{proposition}[Measurable adaptive coordinate charts]
\label{prop:measurable-adaptive-charts}
Under Assumptions~\ref{assump:shared-pfaffian-chain} and
\ref{assump:no-forced-root}, the least finite-speed selector gives a Borel
disjoint partition $K_R=\bigsqcup_{j=1}^NE_j$ with nonzero selected
denominators, the Borel sets $E_{j,m}$ increase exactly to $E_j$, and the
original-coordinate chart \eqref{eq:a4} satisfies \eqref{eq:a5} and
\eqref{eq:a7}.  These conclusions include ties, zero unselected features,
endpoints, empty $K_R$, and $N=1$.
\end{proposition}

\begin{proof}
Lemma~\ref{lem:extended-speed-borel} provides extended-real Borel speeds and
a finite minimum at every feasible point.
Proposition~\ref{prop:lexicographic-pivot-partition} applies the
strict-before and weak-after comparisons in \eqref{eq:a1}, resolving all
ties by the least index.  Lemma~\ref{lem:denominator-exhaustion} proves both
inclusions in \eqref{eq:a3} and finite-level entry.  Finally,
Proposition~\ref{prop:original-coordinate-chart} gives the same-object root
identity, while Proposition~\ref{prop:selected-chart-velocity}
differentiates that expression and proves the literal $R$-weighted bound.
No probability law, area estimate, or chart-count factor is used here.
\end{proof}

\subsection{Measurable Swept Coefficient Volume}

\begin{lemma}[Analytic exhausted charts and root event]
\label{lem:analytic-chart-events}
Under Assumption~\ref{assump:shared-pfaffian-chain} and
Proposition~\ref{prop:measurable-adaptive-charts}, let
$I\subseteq\Theta$ be any interval.  For
\[
J_j=\{1,\ldots,N\}\setminus\{j\},
\qquad B_j=[-R,R]^{J_j},
\]
identify $\mathbb R^{J_j}$ with $\mathbb R^{N-1}$ in increasing original
coordinate order and define
\[
\Psi_j(\theta,\beta)_i=
\begin{cases}
T_j(\theta,\beta),&i=j,\\
\beta_i,&i\in J_j.
\end{cases}
\tag{A8}\label{eq:a8}
\]
For $m\geq1$, set
\[
D_{j,m}=\{(\theta,\beta):
 \theta\in I\cap E_{j,m},\ \beta\in B_j,\
 |T_j(\theta,\beta)|\leq R\}.
\tag{A9}\label{eq:a9}
\]
Then $D_{j,m}$ is Borel in
$\mathbb R\times\mathbb R^{J_j}\cong\mathbb R^N$, and $\Psi_j$ is the
restriction of a $C^1$, locally Lipschitz map on an ambient open
nonzero-pivot chart.  Every $\Psi_j(D_{j,m})$ is analytic and Lebesgue
measurable.  The set
\[
S_I=\{\alpha\in[-R,R]^N:\exists\theta\in I,\
 b(\theta)+\langle\alpha,F(\theta)\rangle=0\}
\tag{A10}\label{eq:a10}
\]
is also analytic and Lebesgue measurable.  This includes open, closed,
half-open, empty, and singleton intervals; endpoints of $\Theta$; cube
faces and corners; empty chart domains; and $N=1$.
\end{lemma}

\begin{proof}
Write $\Theta=[\theta_-,\theta_+]$.  When $\Theta$ is nondegenerate, extend
each $g\in\{b,F_1,\ldots,F_N\}$ to $\mathbb R$ by
\[
\widetilde g(x)=
\begin{cases}
g(\theta_-)+g'(\theta_-)(x-\theta_-),&x<\theta_-,\\
g(x),&\theta_-\leq x\leq\theta_+,\\
g(\theta_+)+g'(\theta_+)(x-\theta_+),&x>\theta_+.
\end{cases}
\tag{A11}\label{eq:a11}
\]
Values and first derivatives agree at both endpoints, so $\widetilde g$ is
$C^1$.  If $\Theta$ is a singleton, use the single affine expression based
at that point and its derivative in the endpoint-inclusive $C^1$
convention.  These extensions agree with the original values and
derivatives throughout $\Theta$.

For each $j$, let
\[
O_j=\{(\theta,\beta)\in
 \mathbb R\times\mathbb R^{J_j}:\widetilde F_j(\theta)\neq0\}.
\tag{A12}\label{eq:a12}
\]
This set is open.  Replacing the functions in \eqref{eq:a4} by their
extensions defines a $C^1$ function $\widetilde T_j$ on $O_j$, and insertion
as in \eqref{eq:a8} defines a $C^1$ map
$\widetilde\Psi_j:O_j\to\mathbb R^N$.  On the selected original chart these
maps equal $T_j$ and $\Psi_j$.  Consequently, even when $\theta$ is an
endpoint, every point of $D_{j,m}$ has an ambient open $C^1$ neighborhood.
This construction imposes no uniform lower bound on a prescribed pivot.

Every interval in $\mathbb R$ is Borel,
Lemma~\ref{lem:denominator-exhaustion} makes $E_{j,m}$ Borel, and $B_j$ is
closed.  Proposition~\ref{prop:lexicographic-pivot-partition} gives
$E_{j,m}\subseteq E_j\subseteq\{F_j\neq0\}$, so the base product in
\eqref{eq:a9} is a Borel subset of $O_j$.  The final condition in
\eqref{eq:a9} is the inverse image of $[-R,R]$ under the continuous ambient
map $\widetilde T_j$.  Thus $D_{j,m}$ is Borel and the restricted map is
Borel.

Kechris's Proposition~14.4 states that if $X$ and $Y$ are Polish spaces,
$f:X\to Y$ is Borel, and $A\subseteq X$ is analytic, then $f(A)$ is
analytic \citep[Section~14, Proposition~14.4, p.~86]{Kechris1995}.  Here
$O_j$ is an open Polish subspace of
$\mathbb R\times\mathbb R^{J_j}$, the Borel set $D_{j,m}\subseteq O_j$ is
analytic in $O_j$, and the global map
$\widetilde\Psi_j:O_j\to\mathbb R^N$ is $C^1$ and hence Borel.  Therefore
$\Psi_j(D_{j,m})=\widetilde\Psi_j(D_{j,m})$ is analytic.

For the full event, define the incidence set
\[
Z_I=\{(\theta,\alpha)\in I\times[-R,R]^N:
 b(\theta)+\langle\alpha,F(\theta)\rangle=0\}.
\tag{A13}\label{eq:a13}
\]
The evaluation map is continuous and the containing product is Borel, so
$Z_I$ is Borel, and hence analytic, in the Polish space
$\Theta\times\mathbb R^N$.  The global coefficient projection
$\pi_\alpha:\Theta\times\mathbb R^N\to\mathbb R^N$ is continuous and
satisfies $\pi_\alpha(Z_I)=S_I$.  The same cited proposition therefore makes
$S_I$ analytic without relying on the graph-coverage identity proved below.
These two applications supply analyticity only; universal and Lebesgue
measurability are obtained separately next, and the coverage, Jacobian, and
volume conclusions are proved later without attribution to Proposition~14.4.

Every set just shown analytic lies in $C=[-R,R]^N$: this is explicit in
\eqref{eq:a10}, and \eqref{eq:a9} puts the beta coordinates and the
inserted $T_j$ coordinate in $C$.  Every analytic subset of a Polish space
is universally measurable
\citep[Section~21, Theorem~21.10, p.~155]{Kechris1995}.  Apply this to the
Borel probability measure
\[
\nu_C(B)=\frac{\lambda_N(B\cap C)}{(2R)^N}.
\]
If $A_0\subseteq C$ is one of the analytic sets, universal measurability
gives a Borel set $B_0$ and a Borel $\nu_C$-null envelope $M$ with
$A_0\mathbin\triangle B_0\subseteq M$.  Since $A_0\subseteq C$,
\[
A_0\mathbin\triangle(B_0\cap C)\subseteq M\cap C,
\qquad \lambda_N(M\cap C)=0.
\]
Hence $A_0$ is Lebesgue measurable in $\mathbb R^N$.

No endpoint or coefficient boundary was removed: the original interval,
closed beta cube, and weak inequality $|T_j|\leq R$ remain in
\eqref{eq:a9}.  For $N=1$, $J_1=\varnothing$ and
$\mathbb R^{J_1}=\mathbb R^0=\{()\}$; the source is still
one-dimensional and the same arguments apply.  Empty domains, images, and
events are analytic and measurable.
\end{proof}

\begin{proposition}[Localized equal-dimensional area sweep]
\label{prop:localized-area-sweep}
Under Assumption~\ref{assump:shared-pfaffian-chain},
Propositions~\ref{prop:original-coordinate-chart} and
\ref{prop:selected-chart-velocity}, and
Lemma~\ref{lem:analytic-chart-events}, order the source coordinates as
\[
\bigl(\theta,\beta_1,\ldots,\beta_{j-1},
\beta_{j+1},\ldots,\beta_N\bigr)
\]
and retain target order $(\alpha_1,\ldots,\alpha_N)$.  For every $j$ and
$m\geq1$,
\[
\int_{D_{j,m}}|\det D\Psi_j(x)|\,d\lambda_N(x)
=\int_{\mathbb R^N}N(\Psi_j,D_{j,m},a)\,d\lambda_N(a),
\tag{A14}\label{eq:a14}
\]
where
\[
N(\Psi_j,D_{j,m},a)
=\#(D_{j,m}\cap\Psi_j^{-1}(a))
\in\{0,1,2,\ldots\}\cup\{+\infty\}.
\]
Moreover,
\[
|\det D\Psi_j(\theta,\beta)|
=|\partial_\theta T_j(\theta,\beta)|,
\tag{A15}\label{eq:a15}
\]
and
\[
\lambda_N(\Psi_j(D_{j,m}))
\leq\int_{D_{j,m}}|\partial_\theta T_j|\,d\lambda_N
\leq(2R)^{N-1}\int_{I\cap E_{j,m}}V_j(\theta)\,d\theta.
\tag{A16}\label{eq:a16}
\]
No injectivity, finite-fiber, regular-value, transversality, or simple-root
condition is imposed.  For $N=1$, \eqref{eq:a15} is scalar and the
zero-dimensional beta measure in \eqref{eq:a16} is one.
\end{proposition}

\begin{proof}
On the open set $O_j$ from Lemma~\ref{lem:analytic-chart-events}, every
target component other than $j$ is its corresponding beta coordinate and
target component $j$ is $\widetilde T_j$.  Move target row $j$ to the first
row, leaving all other target rows in their relative order.  The row
permutation has determinant $(-1)^{j-1}$, and the permuted derivative is
\[
\begin{pmatrix}
\partial_\theta\widetilde T_j&\nabla_\beta\widetilde T_j\\
0&I_{N-1}
\end{pmatrix}.
\]
It has determinant $\partial_\theta\widetilde T_j$, and therefore in the
original target order
\[
\det D\widetilde\Psi_j
=(-1)^{j-1}\partial_\theta\widetilde T_j.
\]
Restriction to the original chart proves \eqref{eq:a15}.  The beta
derivatives occupy only the upper-right block and contribute no determinant
factor.

We next reduce local Lipschitz regularity to the global interface of the
area formula.  Cover $O_j$ by a countable family of open balls
$(\mathcal B_\ell)_{\ell\geq1}$ with rational centers and radii and with
$\overline{\mathcal B_\ell}$ compactly contained in $O_j$.  The derivative
of $\widetilde\Psi_j$ is bounded on each compact closure, so integration of
the derivative along line segments in the convex ball makes
$\widetilde\Psi_j$ Lipschitz there.  Define disjoint Borel layers
\[
\mathcal L_\ell
=\mathcal B_\ell\setminus\bigcup_{r<\ell}\mathcal B_r,
\qquad
D_{j,m}^{(\ell)}=D_{j,m}\cap\mathcal L_\ell.
\tag{A17}\label{eq:a17}
\]
They partition $D_{j,m}$.

For a scalar $L$-Lipschitz coordinate $u$ on $\mathcal B_\ell$, the formula
\[
\widehat u(x)=\inf_{z\in\mathcal B_\ell}
\{u(z)+L|x-z|\}
\tag{A18}\label{eq:a18}
\]
defines a global $L$-Lipschitz extension agreeing with $u$ on the ball.
Applying \eqref{eq:a18} coordinatewise gives a global Lipschitz map
$\widehat\Psi_{j,\ell}:\mathbb R^N\to\mathbb R^N$ that agrees with
$\widetilde\Psi_j$ throughout $\mathcal B_\ell$, and hence has the same
derivative on $D_{j,m}^{(\ell)}$.

The area formula with extended multiplicity, applied to this global map and
the Borel domain piece, gives
\citep[Section~3.2.3, p.~243]{Federer1969}
\[
\int_{D_{j,m}^{(\ell)}}|\det D\Psi_j|\,d\lambda_N
=\int_{\mathbb R^N}
\#(D_{j,m}^{(\ell)}\cap\Psi_j^{-1}(a))\,d\lambda_N(a).
\tag{A19}\label{eq:a19}
\]
Here source and target both have Euclidean dimension $N$; Federer's
normalized $N$-dimensional Hausdorff measure is $\lambda_N$ in this
setting.  Both sides of \eqref{eq:a19} are nonnegative extended integrals.
Summing over $\ell$, exchanging the sum with the nonnegative integral, and
using disjointness of \eqref{eq:a17} yields
\[
\begin{aligned}
\int_{D_{j,m}}|\det D\Psi_j|\,d\lambda_N
&=\int_{\mathbb R^N}\sum_{\ell\geq1}
 \#(D_{j,m}^{(\ell)}\cap\Psi_j^{-1}(a))\,d\lambda_N(a)\\
&=\int_{\mathbb R^N}
 \#(D_{j,m}\cap\Psi_j^{-1}(a))\,d\lambda_N(a),
\end{aligned}
\]
which proves \eqref{eq:a14} with no localization count or constant.

Lemma~\ref{lem:analytic-chart-events} makes the image measurable.  The
multiplicity in \eqref{eq:a14} is at least one on that image, even when it
is infinite.  Thus \eqref{eq:a14} and \eqref{eq:a15} give the first
inequality in \eqref{eq:a16}.  For the second, nonnegative product
integration and Proposition~\ref{prop:selected-chart-velocity} give
\[
\begin{aligned}
\int_{D_{j,m}}|\partial_\theta T_j|\,d\lambda_N
&=\int_{I\cap E_{j,m}}\int_{B_j}
 \mathbf 1_{\{|T_j(\theta,\beta)|\leq R\}}
 |\partial_\theta T_j(\theta,\beta)|
 \,d\lambda_{N-1}(\beta)\,d\theta\\
&\leq\int_{I\cap E_{j,m}}\int_{B_j}V_j(\theta)
 \,d\lambda_{N-1}(\beta)\,d\theta\\
&=\lambda_{N-1}(B_j)
 \int_{I\cap E_{j,m}}V_j(\theta)\,d\theta\\
&=(2R)^{N-1}\int_{I\cap E_{j,m}}V_j(\theta)\,d\theta.
\end{aligned}
\tag{A20}\label{eq:a20}
\]
All faces of $B_j$ and the two faces $T_j=\pm R$ remain in the domain.
For $N=1$, $B_1=[-R,R]^0=\{()\}$ and
$\lambda_0(B_1)=1=(2R)^0$.

For completeness, let $a=\Psi_j(\theta,\beta)$.  Differentiating the root
identity at fixed $\beta$ gives
\[
\partial_\theta T_j(\theta,\beta)
=-\frac{\phi_a'(\theta)}{F_j(\theta)}.
\tag{A21}\label{eq:a21}
\]
Tangent and differentiably multiple roots are therefore critical
preimages.  The Borel critical subset
$D_{j,m}\cap\{\partial_\theta T_j=0\}$ has zero Jacobian, and applying the
already established image inequality to that subset makes its image
$N$-dimensionally null.  If $\phi_a$ vanishes identically on an interval,
the represented fiber may be uncountable.  Finite, countably infinite, and
interval-sized fibers are all permitted by the extended multiplicity in
\eqref{eq:a14}.  Distinct or repeated roots add preimages; the proof never
uses an inverse change of variables.
\end{proof}

\begin{lemma}[Original-coordinate exhausted coverage]
\label{lem:exact-root-event-coverage}
Under Lemma~\ref{lem:root-feasible-compact},
Proposition~\ref{prop:lexicographic-pivot-partition},
Lemma~\ref{lem:denominator-exhaustion}, and
Proposition~\ref{prop:original-coordinate-chart}, for every interval
$I\subseteq\Theta$ and every $j$,
\[
\begin{aligned}
D_{j,m}&\subseteq D_{j,m+1},\\
\Psi_j(D_{j,m})&\subseteq\Psi_j(D_{j,m+1}),
\end{aligned}
\tag{A22}\label{eq:a22}
\]
and
\[
S_I=\bigcup_{j=1}^N\bigcup_{m=1}^\infty\Psi_j(D_{j,m})
\tag{A23}\label{eq:a23}
\]
as literal equality in the original coefficient cube.  This includes
interval and cube boundaries, least-index ties, arbitrarily small nonzero
selected pivots, any number of roots, empty $I$, singleton $I$, empty
$K_R$, empty cells, and $N=1$.
\end{lemma}

\begin{proof}
Only the condition $\theta\in E_{j,m}$ in \eqref{eq:a9} depends on $m$.
Lemma~\ref{lem:denominator-exhaustion} gives
$E_{j,m}\subseteq E_{j,m+1}$; $I,B_j,T_j$, and the weak range condition
remain fixed.  This proves domain nesting, and a fixed map sends nested
domains to nested images.

Take $a\in S_I$.  A witness $\theta\in I$ satisfies $\phi_a(\theta)=0$.
Lemma~\ref{lem:root-feasible-compact} gives $\theta\in K_R$, and
Proposition~\ref{prop:lexicographic-pivot-partition} assigns it to exactly
one least-index cell $E_j$ with $F_j(\theta)\neq0$.  Set
$\beta=a_{-j}\in B_j$ in increasing original-coordinate order.  The root
equation gives
\[
a_j=-\frac{b(\theta)+\sum_{i\neq j}a_iF_i(\theta)}{F_j(\theta)}
=T_j(\theta,\beta).
\tag{A24}\label{eq:a24}
\]
Lemma~\ref{lem:denominator-exhaustion} puts this individual nonzero pivot in
some finite $E_{j,m}$.  Since $|a_j|\leq R$, the pair belongs to
$D_{j,m}$, and \eqref{eq:a8} yields $\Psi_j(\theta,\beta)=a$.  This proves
one inclusion in \eqref{eq:a23}.

Conversely, suppose $a=\Psi_j(\theta,\beta)$ with
$(\theta,\beta)\in D_{j,m}$.  The beta coordinates lie in the closed cube
and $|T_j|\leq R$, so $a\in[-R,R]^N$ coordinate by coordinate, including
faces and corners.  Also $\theta\in I\cap E_{j,m}$, and
Proposition~\ref{prop:original-coordinate-chart} gives
$b(\theta)+\langle a,F(\theta)\rangle=0$.  Thus $a\in S_I$.

The proof chooses an arbitrary existential witness and does not require an
isolated, unique, simple, or regular root.  Images may overlap despite the
parameter-cell tie rule, which does not alter set equality.  A selected
denominator may approach zero, but each nonzero value enters a finite weak
threshold.  Empty $K_R$ makes both sides empty by
Lemma~\ref{lem:root-feasible-compact}; empty or singleton intervals retain
their actual points.  For $N=1$, beta is the empty tuple and
\eqref{eq:a24} is $a_1=-b/F_1$.
\end{proof}

\begin{proposition}[Root-event coefficient-volume certificate]
\label{prop:root-event-volume}
Under Assumptions~\ref{assump:shared-pfaffian-chain} and
\ref{assump:no-forced-root}, Propositions~\ref{prop:root-feasibility-finite-conditioning}
and \ref{prop:measurable-adaptive-charts},
Lemma~\ref{lem:analytic-chart-events},
Proposition~\ref{prop:localized-area-sweep}, and
Lemma~\ref{lem:exact-root-event-coverage}, every interval
$I\subseteq\Theta$ satisfies
\[
\lambda_N(S_I)
\leq(2R)^{N-1}\Gamma_{\rm piv}(b,F;R)|I|.
\tag{A25}\label{eq:a25}
\]
The conclusion retains all boundary, zero-denominator-limit, noninjective,
tangent, multiple-root, infinite-fiber, empty-set, and $N=1$ regimes, with
no chart-count, localization-count, multiplicity, or boundary factor.
\end{proposition}

\begin{proof}
Fix $j$ and set
\[
A_{j,m}=\Psi_j(D_{j,m}),
\qquad A_j=\bigcup_{m=1}^\infty A_{j,m}.
\]
Lemmas~\ref{lem:analytic-chart-events} and
\ref{lem:exact-root-event-coverage} make the $A_{j,m}$ measurable and
increasing.  Continuity of measure from below and
Proposition~\ref{prop:localized-area-sweep} give
\[
\begin{aligned}
\lambda_N(A_j)
&=\lim_{m\to\infty}\lambda_N(A_{j,m})\\
&\leq(2R)^{N-1}\lim_{m\to\infty}
 \int_{I\cap E_{j,m}}V_j(\theta)\,d\theta.
\end{aligned}
\tag{A26}\label{eq:a26}
\]
To avoid any undefined product $0\cdot(+\infty)$, define
\[
h_{j,m}(\theta)=
\begin{cases}
V_j(\theta),&\theta\in I\cap E_{j,m},\\
0,&\theta\notin I\cap E_{j,m}.
\end{cases}
\]
Proposition~\ref{prop:lexicographic-pivot-partition} makes $V_j$ finite on
$E_j$, and Lemma~\ref{lem:denominator-exhaustion} makes the domains increase
to $E_j$.  Thus $h_{j,m}$ are nonnegative Borel functions increasing to the
function equal to $V_j$ on $I\cap E_j$ and zero elsewhere.  Passing to the
limit in their nonnegative integrals turns \eqref{eq:a26} into
\[
\lambda_N(A_j)
\leq(2R)^{N-1}\int_{I\cap E_j}V_j(\theta)\,d\theta.
\tag{A27}\label{eq:a27}
\]
Every denominator boundary is recovered, and $1/m$ leaves no tolerance or
remainder.

By \eqref{eq:a23}, finite subadditivity, and \eqref{eq:a27},
\[
\lambda_N(S_I)
\leq(2R)^{N-1}\sum_{j=1}^N
\int_{I\cap E_j}V_j(\theta)\,d\theta.
\tag{A28}\label{eq:a28}
\]
On $E_j$, Propositions~\ref{prop:lexicographic-pivot-partition} and
\ref{prop:selected-chart-velocity} give
$V_j=\min_iV_i\leq\Gamma_{\rm piv}$.  The Borel cells are disjoint and
cover $K_R$, so
\[
\begin{aligned}
\sum_{j=1}^N\int_{I\cap E_j}V_j(\theta)\,d\theta
&\leq\Gamma_{\rm piv}(b,F;R)\sum_{j=1}^N|I\cap E_j|\\
&=\Gamma_{\rm piv}(b,F;R)|I\cap K_R|\\
&\leq\Gamma_{\rm piv}(b,F;R)|I|.
\end{aligned}
\tag{A29}\label{eq:a29}
\]
Equations \eqref{eq:a28}--\eqref{eq:a29} prove \eqref{eq:a25}.  Chart
images may overlap, but the parameter cells do not; the sum costs one copy
of $|I\cap K_R|$, not $N|I|$.

If $K_R=\varnothing$, Lemma~\ref{lem:root-feasible-compact} gives
$S_I=\varnothing$ and $\Gamma_{\rm piv}=0$.  Empty cells contribute zero.
If $I$ is empty or a singleton, $|I|=0$ and \eqref{eq:a29} makes all chart
integrals zero, so $S_I$ is $N$-dimensionally null.  For $N=1$ there is one
cell and the beta factor is $(2R)^0=1$.  The extended multiplicity already
retains all remaining fiber pathologies.
\end{proof}

\begin{proposition}[Measurable swept coefficient volume]
\label{prop:measurable-swept-volume}
Under Assumptions~\ref{assump:shared-pfaffian-chain} and
\ref{assump:no-forced-root}, for every interval $I\subseteq\Theta$, the
original-coordinate root set $S_I$ in \eqref{eq:a10} is analytic and
Lebesgue measurable and satisfies \eqref{eq:a25}.
\end{proposition}

\begin{proof}
Lemma~\ref{lem:analytic-chart-events} verifies the descriptive-set
interfaces and measurability.  Proposition~\ref{prop:localized-area-sweep}
uses countable disjoint Lipschitz localization to obtain the extended
multiplicity identity, determinant \eqref{eq:a15}, and beta-cube factor.
Lemma~\ref{lem:exact-root-event-coverage} proves both inclusions in
\eqref{eq:a23} in the original coefficient coordinates.
Proposition~\ref{prop:root-event-volume} then performs the increasing limit
and uses the disjoint-cell identity \eqref{eq:a29}.  The separate strict
finiteness of $\Gamma_{\rm piv}$ remains the conclusion of
Proposition~\ref{prop:root-feasibility-finite-conditioning}; the present
result exports measurability and coefficient volume.
\end{proof}

\subsection{Joint-Density Transfer and Ordered Suprema}

\begin{proposition}[Joint-density transfer for the swept root event]
\label{prop:joint-density-transfer}
Under Assumption~\ref{assump:joint-density-cap} and
Proposition~\ref{prop:measurable-swept-volume}, every
$\mu\in\mathcal D_{N,R,\kappa}$ and every interval $I\subseteq\Theta$ with
$|I|>0$ satisfy
\[
\begin{aligned}
\Pr_\mu(S_I)
&\leq\kappa\lambda_N(S_I)\\
&\leq\kappa(2R)^{N-1}\Gamma_{\rm piv}(b,F;R)|I|\\
&=\frac{A\Gamma_{\rm piv}(b,F;R)}{2R}|I|.
\end{aligned}
\tag{A30}\label{eq:a30}
\]
Moreover,
\[
\Pr_{\alpha\sim\mu}\!\left[\exists\theta\in I:
\phi_\alpha(\theta)=0\right]=\Pr_\mu(S_I).
\tag{A31}\label{eq:a31}
\]
These are ordinary probability statements under the completed law.  They
do not require independence, marginal densities, a conditioning event, a
confidence parameter, or a small-interval threshold.
\end{proposition}

\begin{proof}
Fix $\mu$ and $I$, and write $C=[-R,R]^N$ within this proof.
Assumption~\ref{assump:joint-density-cap} gives a nonnegative full joint
density with
\[
f_\mu=0\quad\lambda_N\text{-a.e. on }C^c,
\qquad
f_\mu\leq\kappa\quad\lambda_N\text{-a.e. on }\mathbb R^N.
\tag{A32}\label{eq:a32}
\]
Proposition~\ref{prop:measurable-swept-volume} makes $S_I\subseteq C$
Lebesgue measurable.  Absolute continuity and completion therefore give
\[
\Pr_\mu(S_I)
=\int_{S_I}f_\mu(\alpha)\,d\lambda_N(\alpha)
\leq\kappa\lambda_N(S_I).
\tag{A33}\label{eq:a33}
\]
This uses the cap on the full joint density directly; no coordinate is
integrated out and arbitrary coordinate dependence contributes no term.

The volume certificate in Proposition~\ref{prop:measurable-swept-volume}
applies to the identical set in the identical coefficient coordinates.
Substituting it into \eqref{eq:a33} gives the first two lines of
\eqref{eq:a30}.  Since $R>0$,
\[
\kappa(2R)^{N-1}
=\frac{\kappa(2R)^N}{2R}
=\frac A{2R},
\tag{A34}\label{eq:a34}
\]
which proves the last line without hiding, dropping, or absorbing a term.

Equation \eqref{eq:a32} also gives $\mu(C^c)=0$.  The unrestricted root
event in \eqref{eq:a31} differs from $S_I$ only by a subset of $C^c$; that
subset belongs to the completed law and has probability zero.  This proves
\eqref{eq:a31} without changing the event on the law's support.

If $K_R=\varnothing$, Lemma~\ref{lem:root-feasible-compact} makes $S_I$
empty and the definition gives $\Gamma_{\rm piv}=0$.  If $S_I$ is empty
for another reason, its volume and probability vanish.  If
$\Gamma_{\rm piv}=0$ with nonempty $K_R$, the volume certificate forces
$\lambda_N(S_I)=0$, and \eqref{eq:a33} forces zero probability.  For
$N=1$, $(2R)^{N-1}=1$ and $A/(2R)=\kappa$.  Endpoints, cube faces,
tangent and multiple roots, identically zero affine combinations, and
finite or infinite fibers all remain in $S_I$ and were handled before the
density integration.  The right side of \eqref{eq:a30} may exceed one; it
remains a valid bound and is not clipped.
\end{proof}

\begin{lemma}[Nonemptiness threshold for the law class]
\label{lem:law-class-nonempty}
For $N\geq1$, $R>0$, $\kappa\in(0,\infty)$, and
$A=(2R)^N\kappa$,
\[
\mathcal D_{N,R,\kappa}\neq\varnothing
\quad\Longleftrightarrow\quad A\geq1.
\tag{A35}\label{eq:a35}
\]
The boundary $A=1$ is included, while $A<1$ makes the class empty.
\end{lemma}

\begin{proof}
If $\mu\in\mathcal D_{N,R,\kappa}$ and $C=[-R,R]^N$, then
\[
\begin{aligned}
1
&=\int_{\mathbb R^N}f_\mu(\alpha)\,d\lambda_N(\alpha)\\
&=\int_Cf_\mu(\alpha)\,d\lambda_N(\alpha)\\
&\leq\kappa\lambda_N(C)
=\kappa(2R)^N=A.
\end{aligned}
\tag{A36}\label{eq:a36}
\]
Thus nonemptiness implies $A\geq1$, and correlation cannot alter this
integral constraint.

Conversely, suppose $A\geq1$ and define
\[
f_{\rm unif}(\alpha)=(2R)^{-N}\mathbf 1_{[-R,R]^N}(\alpha).
\tag{A37}\label{eq:a37}
\]
The cube has volume $(2R)^N$, so \eqref{eq:a37} integrates to one and
vanishes outside the cube.  Moreover,
\[
(2R)^{-N}\leq\kappa
\quad\Longleftrightarrow\quad1\leq\kappa(2R)^N=A.
\tag{A38}\label{eq:a38}
\]
Thus it defines an admissible law.  At $A=1$, \eqref{eq:a38} is equality,
proving the sharp boundary without adding $A\geq1$ as an assumption.
\end{proof}

\begin{lemma}[Nonemptiness threshold for positive-length intervals]
\label{lem:interval-class-nonempty}
For a compact interval $\Theta\subset\mathbb R$,
\[
\{I\subseteq\Theta:I\text{ is an interval and }|I|>0\}
\neq\varnothing
\quad\Longleftrightarrow\quad|\Theta|>0.
\tag{A39}\label{eq:a39}
\]
\end{lemma}

\begin{proof}
If the class on the left is nonempty, choose $I$ in it.  Monotonicity of
Lebesgue measure gives $0<|I|\leq|\Theta|$.  Conversely, if
$|\Theta|>0$, the interval $I=\Theta$ belongs to the class.  Hence a
zero-length compact interval has no positive-length subinterval, while a
positive-length compact interval supplies one.  Endpoint and openness
conventions do not affect the equivalence.
\end{proof}

\begin{proposition}[Ordered suprema with nonempty and vacuous semantics]
\label{prop:ordered-suprema}
Under Assumptions~\ref{assump:shared-pfaffian-chain},
\ref{assump:no-forced-root}, and \ref{assump:joint-density-cap}, and under
Propositions~\ref{prop:root-feasibility-finite-conditioning} and
\ref{prop:joint-density-transfer}, the following exhaustive alternatives
hold.  If $A\geq1$ and $|\Theta|>0$, then
\[
0\leq
\sup_{\mu\in\mathcal D_{N,R,\kappa}}
\sup_{\substack{I\subseteq\Theta,\ I\text{ an interval}\\|I|>0}}
\frac{\Pr_{\alpha\sim\mu}[\exists\theta\in I:
\phi_\alpha(\theta)=0]}{|I|}
\leq\frac{A\Gamma_{\rm piv}(b,F;R)}{2R}<\infty.
\tag{A40}\label{eq:a40}
\]
The middle quantity is a finite nonnegative real.  If $A<1$ or
$|\Theta|=0$, then, under $\sup\varnothing=-\infty$,
\[
\sup_{\mu\in\mathcal D_{N,R,\kappa}}
\sup_{\substack{I\subseteq\Theta,\ I\text{ an interval}\\|I|>0}}
\frac{\Pr_{\alpha\sim\mu}[\exists\theta\in I:
\phi_\alpha(\theta)=0]}{|I|}
=-\infty
\leq\frac{A\Gamma_{\rm piv}(b,F;R)}{2R}.
\tag{A41}\label{eq:a41}
\]
The value in \eqref{eq:a41} is a vacuous extended-real consequence, not a
finite nonnegative capacity.  Proposition~\ref{prop:joint-density-transfer}
remains a universal pairwise statement in every branch.
\end{proposition}

\begin{proof}
Suppose first that $A\geq1$ and $|\Theta|>0$.
Lemmas~\ref{lem:law-class-nonempty} and
\ref{lem:interval-class-nonempty} show that both index classes are
nonempty.  For fixed $\mu$, Proposition~\ref{prop:joint-density-transfer}
gives for each indexed interval
\[
0\leq
\frac{\Pr_{\alpha\sim\mu}[\exists\theta\in I:
\phi_\alpha(\theta)=0]}{|I|}
\leq\frac{A\Gamma_{\rm piv}(b,F;R)}{2R},
\tag{A42}\label{eq:a42}
\]
where division is legal because $|I|>0$.  First take the supremum over
intervals for this fixed law:
\[
0\leq
\sup_{\substack{I\subseteq\Theta,\ I\text{ an interval}\\|I|>0}}
\frac{\Pr_{\alpha\sim\mu}[\exists\theta\in I:
\phi_\alpha(\theta)=0]}{|I|}
\leq\frac{A\Gamma_{\rm piv}(b,F;R)}{2R}.
\tag{A43}\label{eq:a43}
\]
Then take the supremum over admissible laws.  This preserves the stated
inner-then-outer order and is not a union bound over either index.
Proposition~\ref{prop:root-feasibility-finite-conditioning} gives
$0\leq\Gamma_{\rm piv}<\infty$ in both feasible-set branches.  Since $A$
is finite and positive and $R>0$, the right side is finite.  Nonempty index
classes and nonnegative ratios make the ordered supremum nonnegative.  If
$\Gamma_{\rm piv}=0$, every ratio is zero and the supremum is zero.

If $A<1$, Lemma~\ref{lem:law-class-nonempty} makes the outer class empty,
so its supremum is $-\infty$, irrespective of the inner class.  If
$|\Theta|=0$, Lemma~\ref{lem:interval-class-nonempty} makes every inner
class empty.  When laws exist, the outer supremum of the constant value
$-\infty$ remains $-\infty$; when the law class is also empty, the outer
empty supremum has the same value.  This covers both single-empty and
double-empty branches.  At $A=1$, the uniform cube law from
\eqref{eq:a37} exists, so the boundary belongs to \eqref{eq:a40} whenever
$|\Theta|>0$.  No branch condition is inserted into the pairwise theorem.
\end{proof}

\begin{proposition}[Density conversion and branch-qualified uniform consequence]
\label{prop:density-and-uniform-consequence}
Under Assumptions~\ref{assump:shared-pfaffian-chain},
\ref{assump:no-forced-root}, and \ref{assump:joint-density-cap}, every
admissible law and positive-length interval satisfy \eqref{eq:a30}--
\eqref{eq:a31}, while the ordered law-then-interval expression has the
semantics in \eqref{eq:a40}--\eqref{eq:a41}, with the interval supremum
taken first.  A finite nonnegative capacity interpretation is valid only
when $A\geq1$ and $|\Theta|>0$.
\end{proposition}

\begin{proof}
Proposition~\ref{prop:measurable-swept-volume} supplies the measurable
original-coordinate event and its coefficient-volume bound.
Proposition~\ref{prop:joint-density-transfer} integrates that same event
against the same capped full joint density and proves the pairwise result
with the algebraic identity \eqref{eq:a34}.  The sharp index classifications
are Lemmas~\ref{lem:law-class-nonempty} and
\ref{lem:interval-class-nonempty}.  Proposition~\ref{prop:ordered-suprema}
then takes the two suprema in their declared order, using strict finiteness
only from Proposition~\ref{prop:root-feasibility-finite-conditioning}.
It retains $-\infty$ whenever an index class is empty and does not reinterpret
that value as a capacity.
\end{proof}

\subsection{Exact Scale-Stress Conditioning}

\begin{lemma}[Root-feasible set for the scale-stress family]
\label{lem:scale-root-feasible}
For $0<\delta\leq1$, let
\[
\Theta=[-1,1],
\qquad b_\delta(\theta)=0,
\qquad F_\delta(\theta)=\left(1,\frac\theta\delta\right),
\qquad R=1.
\]
This family satisfies Assumptions~\ref{assump:shared-pfaffian-chain} and
\ref{assump:no-forced-root}, its feasible set is $K_1=\Theta$, and its
first feature is a legal pivot everywhere on $K_1$.
\end{lemma}

\begin{proof}
For fixed $\delta>0$, the three scalar functions are $C^1$ polynomials in
$\theta$ with a vacuous chain, and the first feature is identically one.
Thus the two stated assumptions hold.  The definition of the feasible set
gives
\[
\begin{aligned}
K_1
&=\{\theta\in\Theta:
 |b_\delta(\theta)|\leq\|F_\delta(\theta)\|_1\}\\
&=\left\{\theta\in\Theta:
 0\leq1+\frac{|\theta|}{\delta}\right\}=\Theta.
\end{aligned}
\tag{A44}\label{eq:a44}
\]
Both inclusions are literal: the inequality holds because $\delta>0$, and
the set comprehension already restricts to $\Theta$.  Hence $K_1$ contains
$0$ and both endpoints.  Since $F_{\delta,1}(\theta)=1$, coordinate one is
nonzero everywhere.  This realizes the legal-pivot conclusion of
Proposition~\ref{prop:root-feasibility-finite-conditioning} without imposing
a global-pivot condition on general families.
\end{proof}

\begin{lemma}[Extended pivot speeds for the scale-stress family]
\label{lem:scale-pivot-speeds}
Under Lemma~\ref{lem:scale-root-feasible}, for every $0<\delta\leq1$,
\[
V_1(\theta)=\frac1\delta\quad(\theta\in\Theta),
\qquad
V_2(\theta)=
\begin{cases}
\displaystyle\frac\delta{\theta^2},&\theta\neq0,\\
+\infty,&\theta=0.
\end{cases}
\tag{A45}\label{eq:a45}
\]
\end{lemma}

\begin{proof}
Write
$F_{\delta,1}(\theta)=1$ and
$F_{\delta,2}(\theta)=\theta/\delta$.  The first coordinate is nonzero on
all of $\Theta$.  With $N=2$ and $R=1$,
\[
V_1(\theta)=
\left|\left(\frac{b_\delta}{F_{\delta,1}}\right)'(\theta)\right|
+\left|\left(\frac{F_{\delta,2}}{F_{\delta,1}}\right)'(\theta)\right|.
\]
Keeping every quotient-rule numerator and denominator visible,
\[
\left(\frac{b_\delta}{F_{\delta,1}}\right)'
=\frac{0\cdot1-0\cdot0}{1^2}=0,
\qquad
\left(\frac{F_{\delta,2}}{F_{\delta,1}}\right)'
=\frac{(1/\delta)\cdot1-(\theta/\delta)\cdot0}{1^2}
=\frac1\delta.
\]
Since $\delta>0$, this proves $V_1=1/\delta$ everywhere, including zero
and both endpoints.

If $\theta\neq0$, then $F_{\delta,2}(\theta)\neq0$ and
\[
V_2(\theta)=
\left|\left(\frac{b_\delta}{F_{\delta,2}}\right)'(\theta)\right|
+\left|\left(\frac{F_{\delta,1}}{F_{\delta,2}}\right)'(\theta)\right|.
\]
The two quotient derivatives are
\[
\left(\frac{b_\delta}{F_{\delta,2}}\right)'
=\frac{0\cdot(\theta/\delta)-0\cdot(1/\delta)}
{(\theta/\delta)^2}=0
\]
and
\[
\begin{aligned}
\left(\frac{F_{\delta,1}}{F_{\delta,2}}\right)'
&=\frac{0\cdot(\theta/\delta)-1\cdot(1/\delta)}
{(\theta/\delta)^2}\\
&=-\frac{1/\delta}{\theta^2/\delta^2}
=-\frac\delta{\theta^2}.
\end{aligned}
\]
Thus $V_2(\theta)=\delta/\theta^2$ off zero.  At zero, the second feature
vanishes, so its quotient is not differentiated; the definition instead
sets $V_2(0)=+\infty$.  Endpoint derivatives use the same $C^1$
convention.
\end{proof}

\begin{proposition}[Exact scale-stress pivot conditioning]
\label{prop:scale-stress-specialization}
Under Proposition~\ref{prop:root-feasibility-finite-conditioning} and
Lemmas~\ref{lem:scale-root-feasible} and \ref{lem:scale-pivot-speeds}, the
following holds for every $0<\delta\leq1$, including $\delta=1$:
\[
\Gamma_{\rm piv}(b_\delta,F_\delta;1)=\frac1\delta.
\tag{A46}\label{eq:a46}
\]
\end{proposition}

\begin{proof}
For every $\theta\in\Theta$,
\[
\min\{V_1(\theta),V_2(\theta)\}
\leq V_1(\theta)=\frac1\delta.
\tag{A47}\label{eq:a47}
\]
At $\theta=0$,
\[
\min\{V_1(0),V_2(0)\}
=\min\left\{\frac1\delta,+\infty\right\}
=\frac1\delta.
\tag{A48}\label{eq:a48}
\]
Lemma~\ref{lem:scale-root-feasible} gives $K_1=\Theta$ and $0\in K_1$.
Consequently,
\[
\frac1\delta
=\min\{V_1(0),V_2(0)\}
\leq\sup_{\theta\in K_1}\min\{V_1(\theta),V_2(\theta)\}
\leq\frac1\delta.
\tag{A49}\label{eq:a49}
\]
The nonempty branch of the definition of $\Gamma_{\rm piv}$ turns
\eqref{eq:a49} into \eqref{eq:a46}.

At zero only pivot one has finite speed.  Away from zero the two speeds tie
exactly when
\[
\frac1\delta=\frac\delta{\theta^2},
\quad\text{that is,}\quad\theta=\pm\delta.
\]
These points lie in $[-1,1]$ for the full range $0<\delta\leq1$; for
$\delta=1$ they are the endpoints.  The least-index convention chooses
pivot one at a tie but does not change the pointwise minimum or supremum.
The coefficient in \eqref{eq:a46} is exactly one.  This is a deterministic
conditioning identity and asserts no additional probability bound.
\end{proof}

\subsection{Exact Affine-Monic Recovery}

\begin{proposition}[Original-space monic presentation]
\label{prop:monic-presentation}
Fix $d\geq1$, $R>0$, $\kappa>0$, and a compact interval
$\Theta\subset\mathbb R$.  Under Assumption~\ref{assump:joint-density-cap},
let $\mu$ be any Borel probability law on $\mathbb R^d$ with full joint
density supported almost everywhere on $[-R,R]^d$ and bounded almost
everywhere by $\kappa$.  Define
\[
p_\alpha(\theta)=\theta^d+\sum_{k=0}^{d-1}\alpha_k\theta^k,
\qquad
b_d(\theta)=\theta^d,
\qquad
F_d(\theta)=(1,\theta,\ldots,\theta^{d-1}).
\tag{A50}\label{eq:a50}
\]
Then for every $\alpha\in\mathbb R^d$ and $\theta\in\Theta$,
\[
b_d(\theta)+\langle\alpha,F_d(\theta)\rangle=p_\alpha(\theta)
\tag{A51}\label{eq:a51}
\]
in the original $d$-dimensional lower-coefficient space.  This family
satisfies Assumptions~\ref{assump:shared-pfaffian-chain} and
\ref{assump:no-forced-root}, and its descriptor tuple is
\[
q=0,
\qquad M=0,
\qquad\Delta_{\rm rnd}=d-1,
\qquad\Delta_{\rm aff}=d,
\qquad N=d,
\qquad A=(2R)^d\kappa.
\tag{A52}\label{eq:a52}
\]
\end{proposition}

\begin{proof}
Index the features from one to $d$, so $(F_d)_i(\theta)=\theta^{i-1}$.
Then
\[
\langle\alpha,F_d(\theta)\rangle
=\sum_{i=1}^d\alpha_{i-1}\theta^{i-1}
=\sum_{k=0}^{d-1}\alpha_k\theta^k,
\]
and adding $b_d(\theta)=\theta^d$ proves \eqref{eq:a51}.  This is equality
for the same $\alpha=(\alpha_0,\ldots,\alpha_{d-1})$.  The leading
coefficient one remains deterministic in $b_d$; there is no additional
coordinate, transformed law, singular augmentation, or coefficient-space
change.

For Assumption~\ref{assump:shared-pfaffian-chain}, take the vacuous-chain
presentation $q=0$, with no $\eta$ and no chain polynomials.  The convention
gives $M=0$, and the output polynomials are
\[
Q_0(\theta)=\theta^d,
\qquad Q_i(\theta)=\theta^{i-1}\quad(1\leq i\leq d).
\]
They are $C^1$ on every compact interval and, in the same scalar convention
$p=1$,
\[
\max_{1\leq i\leq d}\deg Q_i=d-1,
\qquad
\max_{0\leq i\leq d}\deg Q_i=d.
\]
This includes $d=1$, for which the only random feature is the degree-zero
constant.  Hence $\Delta_{\rm rnd}=d-1$, $\Delta_{\rm aff}=d$, and there
are $N=d$ random lower coefficients.

For Assumption~\ref{assump:no-forced-root}, the first feature obeys
$(F_d)_1(\theta)=1$ everywhere.  Hence $F_d(\theta)\neq0$, including at
zero and at both endpoints.  The law hypothesis is the full joint-density
cap in the same coefficient space, with only almost-everywhere cube support.
Substituting $N=d$ into $A=(2R)^N\kappa$ completes \eqref{eq:a52}.  No
coordinate independence has been used or inferred.
\end{proof}

\begin{lemma}[Constant-pivot velocity on the closed unit region]
\label{lem:monic-constant-pivot}
Under Proposition~\ref{prop:monic-presentation}, the constant feature is a
legal pivot at every $\theta$.  If $d=1$, its speed is exactly
\[
V_{\rm const}(\theta)=1.
\tag{A53}\label{eq:a53}
\]
If $d\geq2$ and $|\theta|\leq1$, then
\[
V_{\rm const}(\theta)
=d|\theta|^{d-1}+R\sum_{k=1}^{d-1}k|\theta|^{k-1}
\leq d+R\frac{d(d-1)}2.
\tag{A54}\label{eq:a54}
\]
Thus the displayed upper bound holds for every $d\geq1$ throughout the
closed unit region, with exact speed one for the entire $d=1$ branch.
\end{lemma}

\begin{proof}
The authoritative pivot definition gives
\[
V_{\rm const}(\theta)
=\left|\left(\frac{b_d}{1}\right)'(\theta)\right|
+R\sum_{k=1}^{d-1}
 \left|\left(\frac{\theta^k}{1}\right)'(\theta)\right|.
\tag{A55}\label{eq:a55}
\]
There are exactly $d-1$ nonpivot features; the deterministic leading
monomial appears only in the first term.  For $d=1$, \eqref{eq:a55} is
\[
|\theta'|+R\sum_{k=1}^0(\cdots)=1,
\]
which proves \eqref{eq:a53}.

For $d\geq2$, direct differentiation gives
\[
|(d\theta^{d-1})|=d|\theta|^{d-1},
\qquad
|(k\theta^{k-1})|=k|\theta|^{k-1},
\]
for both signs of $\theta$, proving equality in \eqref{eq:a54}.  At
$\theta=0$ this is read directly: the leading derivative is zero, the
$k=1$ feature derivative is one, and all $k\geq2$ derivatives are zero.
No $0^0$ convention is used.

For $|\theta|\leq1$, every displayed power is at most one, and so
\[
V_{\rm const}(\theta)\leq d+R\sum_{k=1}^{d-1}k.
\]
The arithmetic remains explicit:
\[
2\sum_{k=1}^{d-1}k
=\sum_{k=1}^{d-1}\bigl(k+(d-k)\bigr)
=d(d-1),
\qquad
\sum_{k=1}^{d-1}k=\frac{d(d-1)}2.
\tag{A56}\label{eq:a56}
\]
This proves \eqref{eq:a54}.  Both points $\theta=1$ and $\theta=-1$ stay
in this constant chart; they are not reassigned to the top chart.
\end{proof}

\begin{proposition}[Top-pivot velocity and global monic conditioning]
\label{prop:monic-global-pivot}
Under Proposition~\ref{prop:root-feasibility-finite-conditioning},
Proposition~\ref{prop:monic-presentation}, and
Lemma~\ref{lem:monic-constant-pivot}, if $d\geq2$ and $|\theta|>1$, then
the top feature $(F_d)_d(\theta)=\theta^{d-1}$ is nonzero and
\[
V_{\rm top}(\theta)
=1+R\sum_{m=1}^{d-1}\frac{m}{|\theta|^{m+1}}
\leq1+R\frac{d(d-1)}2
\leq d+R\frac{d(d-1)}2.
\tag{A57}\label{eq:a57}
\]
For every $d\geq1$, $R>0$, and compact interval $\Theta$, the exact
feasible set of the restricted monic family is
\[
K_R=\left\{\theta\in\Theta:
 |\theta|^d\leq R\sum_{k=0}^{d-1}|\theta|^k\right\},
\tag{A58}\label{eq:a58}
\]
and
\[
\Gamma_{\rm piv}(b_d,F_d;R)
\leq d+R\frac{d(d-1)}2.
\tag{A59}\label{eq:a59}
\]
The bound is independent of $\Theta$.
\end{proposition}

\begin{proof}
Assume first $d\geq2$ and $|\theta|>1$.  Then $\theta\neq0$, so the top
feature is a legal pivot.  Its offset ratio is
\[
\frac{b_d(\theta)}{(F_d)_d(\theta)}
=\frac{\theta^d}{\theta^{d-1}}=\theta,
\qquad
\left|\left(\frac{b_d}{(F_d)_d}\right)'(\theta)\right|=1.
\tag{A60}\label{eq:a60}
\]
The lower nonpivot exponents are $k=0,\ldots,d-2$.  Reindex them by
\[
m=d-1-k\in\{1,\ldots,d-1\},
\qquad k-(d-1)=-m.
\tag{A61}\label{eq:a61}
\]
Then
\[
\frac{\theta^k}{\theta^{d-1}}=\theta^{-m},
\qquad
| (\theta^{-m})' |
=|-m\theta^{-m-1}|
=\frac{m}{|\theta|^{m+1}}.
\tag{A62}\label{eq:a62}
\]
Substituting \eqref{eq:a60}--\eqref{eq:a62} into the definition of $V_j$
gives equality in \eqref{eq:a57} for both signs of $\theta$.  Since
$|\theta|>1$, each $m/|\theta|^{m+1}\leq m$; equation \eqref{eq:a56}
gives the first inequality, and $1\leq d$ gives the second.  There is no
remainder or large-radius threshold.

Substitution of the monic objects into the definition of $K_R$ gives
\[
|b_d(\theta)|=|\theta|^d,
\qquad
\|F_d(\theta)\|_1
=\sum_{k=0}^{d-1}|\theta^k|
=\sum_{k=0}^{d-1}|\theta|^k,
\]
which is \eqref{eq:a58}.  At every point of this set take the minimum over
legal pivots.  When $d=1$, the constant feature is the sole feature and
Lemma~\ref{lem:monic-constant-pivot} gives
\[
\min_{j:(F_1)_j(\theta)\neq0}V_j(\theta)
=1=d+R\frac{d(d-1)}2.
\tag{A63}\label{eq:a63}
\]
When $d\geq2$ and $|\theta|\leq1$, the constant feature is legal and
\[
\min_{j:(F_d)_j(\theta)\neq0}V_j(\theta)
\leq V_{\rm const}(\theta)
\leq d+R\frac{d(d-1)}2.
\tag{A64}\label{eq:a64}
\]
When $d\geq2$ and $|\theta|>1$, the top feature is legal and
\[
\min_{j:(F_d)_j(\theta)\neq0}V_j(\theta)
\leq V_{\rm top}(\theta)
\leq d+R\frac{d(d-1)}2.
\tag{A65}\label{eq:a65}
\]
At $|\theta|=1$, only \eqref{eq:a64} is used, so there is no duplicated
boundary or chart-count cost.

If $K_R$ is nonempty, taking its defining supremum in
\eqref{eq:a63}--\eqref{eq:a65} proves \eqref{eq:a59}.  If it is empty,
the definition and Proposition~\ref{prop:root-feasibility-finite-conditioning}
give $\Gamma_{\rm piv}=0$.  A chart region may itself be empty and then
contributes nothing.

For $d=2$, the two branches reduce explicitly to
\[
V_{\rm const}(\theta)=2|\theta|+R\leq2+R
\quad(|\theta|\leq1),
\]
\[
V_{\rm top}(\theta)=1+\frac R{|\theta|^2}
\leq1+R\leq2+R
\quad(|\theta|>1).
\]
Thus the first nontrivial dimension is covered.  The right side of
\eqref{eq:a59} depends only on $d$ and $R$, not the location, diameter, or
endpoints of $\Theta$.
\end{proof}

\begin{proposition}[Positive-length monic probability transfer]
\label{prop:monic-positive-interval}
Under Assumption~\ref{assump:joint-density-cap},
Propositions~\ref{prop:monic-presentation}, \ref{prop:monic-global-pivot},
and \ref{prop:joint-density-transfer}, let $I\subset\mathbb R$ be a bounded
interval with $|I|>0$.  For the same arbitrary full joint law $\mu$ on the
original $d$ lower coefficients,
\[
\Pr_{\alpha\sim\mu}\!\left[\exists\theta\in I:p_\alpha(\theta)=0\right]
\leq\kappa(2R)^{d-1}
\left(d+R\frac{d(d-1)}2\right)|I|.
\tag{A66}\label{eq:a66}
\]
Equivalently,
\[
\kappa(2R)^{d-1}
\left(d+R\frac{d(d-1)}2\right)|I|
=\frac A{2R}
\left(d+R\frac{d(d-1)}2\right)|I|.
\tag{A67}\label{eq:a67}
\]
\end{proposition}

\begin{proof}
Set $a=\inf I$ and $c=\sup I$.  Boundedness makes both finite.  If $a=c$,
the between-points property of an interval would make $I$ contain at most
one point and hence have length zero, contrary to the hypothesis.  Thus
$a<c$, and the interval property gives $\overline I=[a,c]$.  Use the
nondegenerate compact localization $\Theta=[a,c]$.

Every hypothesis of Proposition~\ref{prop:joint-density-transfer} now holds
for one fixed admissible pair.  Proposition~\ref{prop:monic-presentation}
gives $N=d\geq1$, the $C^1$ vacuous-chain presentation on this $\Theta$,
and no-forced-root nondegeneracy from the constant feature.  The given law
has the capped full joint density in the same original coefficient space,
with almost-everywhere support on $[-R,R]^d$ and no independence premise.
The interval $I$ lies in $\Theta$ and has positive length.
Proposition~\ref{prop:monic-global-pivot} applies to this restriction and
gives
\[
\Gamma_{\rm piv}(b_d,F_d;R)
\leq d+R\frac{d(d-1)}2,
\]
independently of the chosen localization.  Finally, \eqref{eq:a51}
identifies the affine event with the polynomial event for the same
$\alpha,\theta,I$, and $\mu$.

Applying \eqref{eq:a30} with $N=d$ therefore gives
\[
\begin{aligned}
\Pr_{\alpha\sim\mu}[\exists\theta\in I:p_\alpha(\theta)=0]
&=\Pr_{\alpha\sim\mu}[\exists\theta\in I:
 b_d(\theta)+\langle\alpha,F_d(\theta)\rangle=0]\\
&\leq\kappa(2R)^{d-1}
 \Gamma_{\rm piv}(b_d,F_d;R)|I|\\
&\leq\kappa(2R)^{d-1}
 \left(d+R\frac{d(d-1)}2\right)|I|,
\end{aligned}
\]
which is \eqref{eq:a66}.  Because $R>0$,
\[
\kappa(2R)^{d-1}
=\frac{\kappa(2R)^d}{2R}=\frac A{2R},
\]
proving \eqref{eq:a67} while preserving the requested $\kappa$ form.

If $I\cap K_R=\varnothing$, Lemma~\ref{lem:root-feasible-compact} rules out
roots on $I$ for every coefficient in the cube.  Almost-everywhere cube
support assigns probability one to that cube, so the unrestricted event has
probability zero.  This includes empty $K_R$ and shows that no nonemptiness
condition entered the argument.  The right side of \eqref{eq:a66} may
exceed one; no clipping, small-interval condition, or probability-mode
change is made.  The ordered supremum and its capacity interpretation are
not used in this specialization.
\end{proof}

\begin{proposition}[Complete exact monic baseline]
\label{prop:complete-monic-baseline}
Under Assumption~\ref{assump:joint-density-cap},
Proposition~\ref{prop:monic-presentation}, and
Proposition~\ref{prop:monic-positive-interval}, every bounded interval
$I\subset\mathbb R$, including $|I|=0$, satisfies
\[
\Pr_{\alpha\sim\mu}\!\left[\exists\theta\in I:p_\alpha(\theta)=0\right]
\leq\kappa(2R)^{d-1}
\left(d+R\frac{d(d-1)}2\right)|I|.
\tag{A68}\label{eq:a68}
\]
For a singleton interval, its coefficient root set is a proper affine
hyperplane of $\mathbb R^d$ and has both $d$-dimensional Lebesgue measure
and $\mu$-probability zero, including $d=1$.
\end{proposition}

\begin{proof}
Positive-length intervals are covered by
Proposition~\ref{prop:monic-positive-interval}.  It remains to treat zero
length.  If a real interval contained two distinct points $x<y$, it would
contain $(x,y)$ and have length at least $y-x>0$.  Hence a zero-length
interval is empty or a singleton.  The root event is empty when
$I=\varnothing$.

Suppose $I=\{\theta_0\}$.  Equation \eqref{eq:a51} gives the coefficient
root set
\[
\begin{aligned}
H_{\theta_0}
&=\{\alpha\in\mathbb R^d:p_\alpha(\theta_0)=0\}\\
&=\left\{\alpha\in\mathbb R^d:
 \alpha_0=-\theta_0^d-\sum_{k=1}^{d-1}\alpha_k\theta_0^k\right\}.
\end{aligned}
\tag{A69}\label{eq:a69}
\]
The coefficient of $\alpha_0$ is one.  Equivalently, this is an affine
level set with nonzero normal
$F_d(\theta_0)=(1,\theta_0,\ldots,\theta_0^{d-1})$.  It is therefore
proper for every $\theta_0$, including zero and $d=1$.

We prove nullity in these coordinates.  If $d=1$, then
$H_{\theta_0}=\{-\theta_0\}$.  For every $\varepsilon>0$, this point is
covered by an interval of length $\varepsilon$, so its outer measure is at
most $\varepsilon$; letting $\varepsilon\downarrow0$ gives
$\lambda_1(H_{\theta_0})=0$.

If $d\geq2$, write
$\beta=(\alpha_1,\ldots,\alpha_{d-1})\in\mathbb R^{d-1}$.  The set in
\eqref{eq:a69} is the graph of the continuous affine function
\[
\beta\longmapsto-\theta_0^d-
\sum_{k=1}^{d-1}\beta_k\theta_0^k.
\]
It is closed and Lebesgue measurable.  For fixed $\beta$, its
$\alpha_0$-section is a singleton and has one-dimensional measure zero.
Nonnegative product integration of its indicator yields
\[
\begin{aligned}
\lambda_d(H_{\theta_0})
&=\int_{\mathbb R^{d-1}}\int_{\mathbb R}
 \mathbf 1_{H_{\theta_0}}(\alpha_0,\beta)
 \,d\alpha_0\,d\beta\\
&=\int_{\mathbb R^{d-1}}0\,d\beta=0.
\end{aligned}
\tag{A70}\label{eq:a70}
\]
Thus nullity is derived rather than assumed.

The full joint density now gives
\[
\Pr_{\alpha\sim\mu}[p_\alpha(\theta_0)=0]
=\mu(H_{\theta_0})
=\int_{H_{\theta_0}}f_\mu\,d\lambda_d=0.
\tag{A71}\label{eq:a71}
\]
No marginal density or independence is used.  Almost-everywhere cube
support creates no exceptional mass because the law is absolutely
continuous.  The right side of \eqref{eq:a68} is also zero when
$|I|=0$, completing the empty and singleton branches.
\end{proof}

\begin{proposition}[Exact affine-monic specialization and baseline bridge]
\label{prop:exact-monic-recovery}
Under Assumption~\ref{assump:joint-density-cap},
Proposition~\ref{prop:root-feasibility-finite-conditioning}, and the
pairwise probability result in Proposition~\ref{prop:joint-density-transfer},
the monic family in \eqref{eq:a50} has the original-space identity and
descriptors \eqref{eq:a51}--\eqref{eq:a52}, obeys the conditioning bound
\eqref{eq:a59} for every compact restriction, and satisfies the all-bounded-
interval probability bound \eqref{eq:a68}.  The coefficient is literally
\[
\kappa(2R)^{d-1}\left(d+R\frac{d(d-1)}2\right),
\]
with no hidden constant, transformed probability, random leading
coefficient, independence condition, or conservative remainder.
\end{proposition}

\begin{proof}
Proposition~\ref{prop:monic-presentation} keeps the leading coefficient in
$b_d$, keeps exactly $d$ lower coefficients in the random vector, proves
$b_d+\langle\alpha,F_d\rangle=p_\alpha$, checks both deterministic
assumptions, and computes the descriptor tuple.  Lemma~\ref{lem:monic-constant-pivot}
computes the closed-unit-region speed, including $d=1$, zero, both signs,
and $|\theta|=1$.  Proposition~\ref{prop:monic-global-pivot} computes the
top-pivot speed with the reindexing \eqref{eq:a61}, takes the pointwise
minimum over legal pivots, and then takes the supremum over the exact
$K_R$, proving \eqref{eq:a59} with no localization dependence.

For positive-length bounded intervals,
Proposition~\ref{prop:monic-positive-interval} checks the compact
localization and every premise of the pairwise probability result, then
uses the same objects and same law to obtain \eqref{eq:a66} and the exact
algebra \eqref{eq:a67}.  Proposition~\ref{prop:complete-monic-baseline}
adds empty and singleton intervals by the explicit nullity calculation
\eqref{eq:a69}--\eqref{eq:a71}.  No ordered-supremum or capacity branch is
used.  The bridge has ordinary probability for each arbitrary full joint
law, static uniformity over bounded intervals, the scalar pivot-variation
and full $d$-dimensional Lebesgue modes, and no auxiliary parameter or
hidden constant.

This fixed-family bridge does not give polynomial general-instance control
of $\Gamma_{\rm piv}$ in the descriptors
\[
q,\quad M,\quad\Delta_{\rm rnd},\quad\Delta_{\rm aff},
\]
or in any other general Pfaffian presentation data.  That source-direction
gap remains unresolved.
\end{proof}

\subsection{Proof of the Main Theorem}

\begin{proof}[Proof of Theorem~\ref{thm:main}]
Proposition~\ref{prop:root-feasibility-finite-conditioning} proves, from
Assumptions~\ref{assump:shared-pfaffian-chain} and
\ref{assump:no-forced-root}, that every cube-supported root is feasible.  It
also proves, including the empty-$K_R$ branch, that
\[
0\leq\Gamma_{\rm piv}(b,F;R)<\infty.
\]
This is the first assertion of Theorem~\ref{thm:main}.

Proposition~\ref{prop:measurable-adaptive-charts} constructs the disjoint
Borel pivot cells, their complete finite-denominator exhaustion, and the
original-coordinate graph maps from that finite-conditioning interface.
Using precisely those charts, Proposition~\ref{prop:measurable-swept-volume}
proves measurability of the root event and its coefficient-volume bound.
Proposition~\ref{prop:density-and-uniform-consequence} combines that volume
result with the primitive full joint-density cap, while using
Proposition~\ref{prop:root-feasibility-finite-conditioning} directly for
strict finiteness.  Its pairwise statement is
\[
\Pr_{\alpha\sim\mu}[\exists\theta\in I:\phi_\alpha(\theta)=0]
\leq\kappa(2R)^{N-1}\Gamma_{\rm piv}(b,F;R)|I|
=\frac{A\Gamma_{\rm piv}(b,F;R)}{2R}|I|,
\]
which is \eqref{eq:t1}.  Its sharp law- and interval-index classifications
then give \eqref{eq:t2} with a finite nonnegative value exactly when
$A\geq1$ and $|\Theta|>0$, and with the literal value $-\infty$ in every
empty-index branch.

Proposition~\ref{prop:scale-stress-specialization}, based on the common
finite-conditioning result and the two explicit specialized pivot speeds,
proves \eqref{eq:t3} for the entire range $0<\delta\leq1$.
Proposition~\ref{prop:exact-monic-recovery} uses the same common
root-feasibility result and only the pairwise probability statement of
Proposition~\ref{prop:joint-density-transfer}.  It proves the two-pivot
bound \eqref{eq:t4}, preserves the deterministic leading coefficient, and
derives \eqref{eq:t5} for positive-length intervals before completing empty
and singleton intervals by full-joint-density nullity.

All coefficients in these cited propositions are displayed literally:
there is no hidden constant, probability-mode conversion, horizon upgrade,
independence premise, or suppressed norm change.  Their fixed-family scope
also leaves the stated absence of a polynomial general Pfaffian-format
bound intact.  This proves every clause of Theorem~\ref{thm:main}.
\end{proof}
